% ============================================================================
% Part 12: Design critique (harsh but constructive)
% ============================================================================
\chapter{Design critique: what is wrong here, and how to fix it}
\label{ch:critique}

The preceding chapters establish that the generated verifier is
\emph{correct} --- no bug was found that breaks soundness on the
reviewed artifacts. This chapter is not about correctness. It is
about whether the \textbf{design} is sound, and it pulls no
punches. The thesis: \emph{the architecture optimizes for the
wrong axis (contract size) and pays for it in auditability,
runtime cost, and unverified trust roots --- and it papers over
each cost with a band-aid instead of fixing the cause.}

A correct verifier that nobody can audit, that no tool can
analyze, whose soundness level under truncated challenges is unstated, and whose
trust root is a hand-written interpreter with no formal proof, is
a liability even when it has no bug today. Each critique below is
tagged \textcolor{rustorange}{\textbf{[CRITICAL]}} /
\textcolor{mathpurple}{\textbf{[MAJOR]}} /
\textcolor{codegenblue}{\textbf{[MINOR]}} and paired with a
concrete, actionable fix.

\paragraph{Status tags (re-verified).} Every item and fix also
carries a status tag checked against generator \texttt{main} HEAD
\texttt{6b3d2096} (2026-07-17) and the B64 render regenerated
2026-07-28: \vstatus{ADDRESSED} (the current code already does
it), \vstatus{PARTIAL}, \vstatus{OPEN}, or \vstatus{CORRECTED}
(the critique's own factual premise was wrong and has been
rewritten). B64 was rendered from the M7 follow-up line
(\texttt{e7f736c4}), not from \texttt{main}; the mechanisms cited
below exist in both, except the quotient-certificate-hash check
(B64 L940), which only the render lineage has. \texttt{file:line}
references are to \texttt{main} HEAD. Line numbers \texttt{L}$n$ refer to
\texttt{Halo2VerifierFee12345B64.sol}: 3{,}579 lines (\texttt{wc -l}
prints 3{,}578 because the closing brace has no trailing newline);
\texttt{verifyProof} spans L293--L3578, i.e.\ \textbf{3{,}286 lines}.

% ---------------------------------------------------------------------------
\section{C1 --- The VM is the wrong abstraction}
\label{sec:c1-vm}
\textcolor{mathpurple}{\textbf{[MAJOR]}} \vstatus{PARTIAL} \emph{(downgraded from
CRITICAL --- see the severity reassessment below)}

The ``generated verifier'' is not a verifier. It is a
\textbf{bytecode interpreter} --- the VM dispatch loop is
$\approx$873 lines of Yul (L1692--2564): $\approx$270 lines of
fixed interpreter template (L1692--1929, L2531--2564) plus
$\approx$600 lines of circuit-specific native-callback Yul emitted
by the generator into the \texttt{0x19}/\texttt{0x1f}/\texttt{0x1b}
arms (L1930--2530) --- plus an
opaque \texttt{q\_program} blob of \textbf{4{,}559 bytes stored as
143 hex words in the VK contract}. The actual verification logic
is \textbf{invisible in the contract}: you cannot read what the
verifier checks by reading the contract. You must disassemble
4{,}559 bytes of bytecode against the VM ISA (31 opcodes defined,
\texttt{vm/\allowbreak{}mod.rs:560--590}; only the 11 the program uses are
rendered in B64) and mentally
execute the VM to recover the semantics. \textbf{This is the
real cost of the design, and it is large.}

\paragraph{The runtime-cost half of the critique is FALSE (measured)} \vstatus{CORRECTED}
My original framing claimed the VM is ``paid for at runtime.''
Measurement refutes the significance of that claim, though the
numbers first given here were wrong and are replaced:
\begin{itemize}[leftmargin=1.4em,itemsep=2pt]
\item Decoding the VK's \texttt{q\_program} with the operand widths
  of the rendered arms ends exactly at byte 4{,}559 and yields only
  \textbf{101 instructions} (not 1{,}139): 16 \texttt{MODARITH7}
  macro-instructions carry 4{,}331\,B (95\%) of limb-term operands
  (1{,}151 terms), the rest are 85 short ops and native-callback
  markers.
\item Dispatch cost $\approx$ \texttt{mload} + \texttt{byte} + an
  \textbf{11-arm} switch per instruction (not 31-way). Even if every
  dispatch walked all 11 compare-and-branch steps
  ($\approx$20 gas each), that is $\lesssim$225 gas per instruction,
  i.e.\ $\lesssim$23k gas over 101 dispatches.
\item Precompile gas $\approx$ \textbf{709.6k} (EIP-2537/EIP-2565
  prices as in \texttt{revm-precompile} 16.2.0): 78-pair G1MSM
  525{,}096 + six 1-pair G1MSMs (L827, 872, 3475, 3486, 3531, 3545)
  72{,}000 + \textbf{one} 2-pair
  pairing call 102{,}900 (the accumulator equation is folded into
  it by a random $\alpha$, L3497--3553; the other \texttt{0x0f} call
  is the constructor smoke test, L250) + four G1ADDs 1{,}500 + six
  modexp inversions $\approx$8{,}124 (24{,}384 under EIP-7883, now live,
  for a total of $\approx$725.9k).
\item Measured \texttt{verifyProof} gas for B64 is $\approx$1.14M
  (\texttt{nightfall\_\allowbreak{}4\_\allowbreak{}PV/\allowbreak{}audit/\allowbreak{}h2v-fee12345/\allowbreak{}evidence/\allowbreak{}bundle/\allowbreak{}logs-rerun/\allowbreak{}p1\_\allowbreak{}01\_\allowbreak{}abi.log}:
  1{,}140{,}607). \textbf{Dispatch is $\lesssim$2\% of total}; precompiles
  are $\approx$62\%, not ``40$\times$'' the rest.
\end{itemize}
So the \emph{dispatch} tax is negligible. The interpretive cost
that remains is inside \texttt{MODARITH7}, which decodes its 1{,}151
terms at run time (an \texttt{mload}/\texttt{byte}/\texttt{shr}
per term); it is unmeasured for B64 but bounded well below the
precompile floor. A critique that leans on ``the VM is slow'' is
still wrong and should not be made.

\paragraph{The generality-unused point} \vstatus{CORRECTED}
The \texttt{EXPECTED\_\allowbreak{}VK\_\allowbreak{}CODEHASH\_\allowbreak{}WORD} anchor pins
\textbf{one VK per verifier deployment} (\S\ref{sec:pr-vk-generic}),
and the verifier also hard-codes circuit shape (e.g.\
\texttt{let k := 21} at L1322 and L2846; the pruned 11-arm switch).
But the VM was never meant to serve many circuits: the header
says ``It is not a generic verifier'' (L11--13). Its justification
is size, because the program lives in VK \emph{data} rather than
verifier \emph{code}. So ``unused generality'' is not a design
smell here. What remains is the cost below: the program is opaque.

\paragraph{Steelman (stronger than I first credited).}
The VM is justified by the 24\,KB EIP-170 limit: a fully
specialized verifier for this 15-advice-column, degree-5 circuit
with all constraints inlined would very likely not fit (B64's
runtime is already 21{,}253\,B). Combined with the
$\lesssim$2\% dispatch overhead, the VM is a
\textbf{reasonable engineering choice on size and cost}. The
failure is not the VM itself --- it is that \textbf{the size win
was paid for in auditability and never bought back}: the program
blob is not human-auditable. A size-driven VM is fine \emph{if} you
then make the program legible --- which this design does not.

\paragraph{Fix (refocused on the real cost: opacity, not speed).}
\begin{itemize}[leftmargin=1.4em,itemsep=2pt]
\item \textbf{Emit an annotated disassembly appendix.} The
  generator knows every opcode's meaning; emit the
  \texttt{q\_program} as a commented block (in the VK or a
  companion file), each instruction tagged with the circuit
  constraint it enforces. This directly attacks the opacity at
  zero runtime cost --- the single highest-value C1 fix.
  \vstatus{PARTIAL}: the host-side
  \texttt{quotient\_\allowbreak{}identity\_\allowbreak{}manifest()} API
  (\texttt{builder/\allowbreak{}api.rs:121}) already maps every $y$-batch
  position to its gate/constraint name, and the canonical
  quotient-certificate hash (identity order, VM bytecode, constant
  table) is checked on-chain against the VK header (L940); what is
  missing is an opcode-level listing shipped with the artifact.
\item \textbf{Formally verify the interpreter} --- it is the trust
  root and the opcode set is bounded. Note the scope: the bounded
  ISA is the $\approx$270-line template core; the $\approx$600
  lines of native-callback Yul are regenerated per circuit and need
  per-artifact checking (differential or translation validation),
  not a one-off ISA proof. \vstatus{OPEN}
\item \textbf{Differentially test the shipped program}
  (\S\ref{sec:c-scorecard}): the VM is data-oblivious, so one passing
  trace run of B64's own program shows its whole
  disassembly is exercised. \vstatus{OPEN}
\item \textbf{Document the size budget} in the contract header so
  a reviewer understands the VM is a 24\,KB necessity, not
  gratuitous indirection. \vstatus{OPEN}
\item \textbf{Reconsider the per-VK pinning} if multi-VK reuse
  was ever intended --- otherwise drop the pretense of generality
  and document the verifier as single-VK. \vstatus{ADDRESSED}: the
  NatSpec already says ``It is not a generic verifier \ldots\
  generated for one \texttt{VerifyingKey}'' (L11--13).
\end{itemize}

\paragraph{Severity reassessment (why MAJOR, not CRITICAL).}
A CRITICAL rating should mean \emph{a flaw that, left in place,
accepts invalid proofs or loses funds}. C1 is not that. The
verifier is \textbf{correct} --- the differential harness passes
on the generator's fixtures and B64 passes Nightfall's own proof
fixtures. The VM is \textbf{size-justified} (24\,KB limit) and
\textbf{cheap} ($\lesssim$2\% dispatch). What C1 actually is: a
\textbf{high auditability risk that amplifies the genuinely
critical issues}. The opaque program blob and the 3{,}286-line
\texttt{verifyProof} do not break soundness --- but they make
the real soundness risks (C3 truncated challenges with an unstated level, C4
unenforced invariants) \emph{harder for a human to spot}. A risk
multiplier is serious, but it is not itself the fire. Rating it
CRITICAL conflated ``hard to audit'' with ``broken,'' which is
the kind of imprecision that erodes a critique's credibility.
MAJOR is the honest grade: fix it because it makes everything
else auditable, not because it is on fire.

\paragraph{Remediation plan (prioritized, with effort and impact).}
\begin{center}\small
\begin{tabular}{@{}cp{4.3cm}cp{3.3cm}p{1.9cm}@{}}
\toprule
\textbf{\#} & \textbf{Action} & \textbf{Effort} & \textbf{Buys} & \textbf{Status} \\
\midrule
R1 & Emit an \textbf{annotated disassembly} of \texttt{q\_program}
  (opcode-by-opcode, each tagged with the circuit constraint) as
  a companion file or VK comment block & M & Kills the opacity:
  a reviewer reads constraints, not hex & \vstatus{PARTIAL} \\
R2 & \textbf{Per-artifact differential run}: trace-compare each
  shipped program (B64 included) against Rust in CI; a census of its
  \texttt{used\_ops} (\texttt{vm/mod.rs:1217}; 11 in B64) shows what it
  covers & S &
  Turns ``tested some shapes'' into ``tested the program we ship'' & \vstatus{OPEN} \\
R3 & \textbf{Split \texttt{verifyProof}} (3{,}286 lines) into
  named sub-functions (transcript, quotient, pairing, acc) & M &
  Auditable units; may \emph{reduce} (not remove) the via-IR
  dependency, see C11 & \vstatus{OPEN} \\
R4 & \textbf{Formal model of the VM ISA} ($\approx$270-line
  template core, bounded opcode set) & L & Proves the trust root, not just
  tests it & \vstatus{OPEN} \\
R5 & \textbf{Document the size budget} + single-VK framing in
  the header & S & Sets the reviewer's mental model right & \vstatus{PARTIAL} \\
\bottomrule
\end{tabular}
\end{center}
\textbf{Do R1 and R2 first} --- they are small, attack the real
cost (opacity + unverified coverage), and together make the
contract cross-checkable against the Rust verifier by human eyes.
R3 is the structural fix; it helps, but does not by itself
remove, the build fragility (C11).
R4 is the gold-standard close for the trust root. None of these
change the bytecode semantics --- they change what the generator
\emph{emits}, which is the whole point: the artifact can be made
legible without touching what it computes.

% ---------------------------------------------------------------------------
\section{C2 --- The trust root is differentially tested, but not formally verified}
\label{sec:c2-trustroot}
\textcolor{mathpurple}{\textbf{[MAJOR]}} \vstatus{PARTIAL} \emph{(revised downward from
CRITICAL after inspecting the CI)}

Every proof this contract accepts is correct \emph{only if the
Yul interpreter (template core plus generated callbacks) is correct}. The VM is the actual
trust root. \textbf{Credit where due:} the generator's test
harness, as wired in
\texttt{proofs/\allowbreak{}solidity-verifier/\allowbreak{}.github/\allowbreak{}workflows/\allowbreak{}ci.yml}, is
genuinely rigorous --- it runs a property-based differential
fuzzer (\texttt{pbt\_\allowbreak{}transcript\_\allowbreak{}differential\_\allowbreak{}fuzzer\_\allowbreak{}matches\_\allowbreak{}solidity\_\allowbreak{}trace})
with \texttt{HALO2\_\allowbreak{}SOLIDITY\_\allowbreak{}RUN\_\allowbreak{}EVM\_\allowbreak{}TESTS=1}, plus dedicated
\emph{native-vs-Solidity trace-equivalence} jobs for Poseidon and
IVC compiled with
\texttt{rust-verifier-trace,solidity-trace}. The comparison is
bidirectional: every Rust trace id must appear in the Solidity
trace, the Solidity trace must emit no id without a Rust
oracle, and every shared id must carry equal bytes
(\texttt{src/\allowbreak{}test.rs:1060--1103}). That is a high bar.
\emph{Correction:} it does \emph{not} run on every PR. GitHub executes
only workflows under the repository root's \texttt{.github/workflows},
the nested \texttt{ci.yml} is inert in this monorepo, and the root
\texttt{ci.yaml} has no solidity-verifier job. The harness exists and
passes locally, but nothing gates PRs on it.

\paragraph{The residual gap (why this is still a critique).}
Differential testing proves equivalence \emph{on the shapes and
traces exercised}, not a formal proof for all inputs. Two
specific holes:
\begin{itemize}[leftmargin=1.4em,itemsep=2pt]
\item \textbf{Program-scoped coverage, not opcode-scoped.} The
  generator renders only the switch arms a program uses
  (\texttt{quotient.rs:194--229}), so B64 contains 11 of the
  31 opcodes and every other byte hits the default revert (L2561);
  an unrendered opcode cannot execute. The VM also has no jumps, and
  every branch in the dispatch body depends only on program bytes,
  so one honest run of a program executes \emph{all} of its rendered
  arms and MODARITH7 sub-blocks. The question is therefore not
  ``untested opcodes'' but ``has \emph{this} program been
  differentially tested?''. The closest fixture, the IVC Keccak
  decider (\texttt{tests/\allowbreak{}ivc\_\allowbreak{}keccak\_\allowbreak{}solidity.rs}, $k{=}20$), renders
  a 4{,}559-byte program with the same 11-opcode histogram as B64
  (checked against the local render in
  \texttt{target/\allowbreak{}ivc-keccak-solidity-dump}). It is not byte-identical,
  though: B64's $k{=}21$, 19-instance layout, embedded-quotient mode and
  exact generator revision are never differentially tested, and the
  other fixtures (Poseidon $k{=}6$, RSA/SHA/hybrid-MT $k{=}12$--13) are
  far smaller.
\item \textbf{Direct coverage is against a model, not the Yul.}
  \texttt{vm/mod.rs} itself has 2 unit tests, but
  \texttt{src/\allowbreak{}lowering/\allowbreak{}tests.rs} has over 20 VM tests,
  including per-opcode encoder tests for LIN7, BILIN7\_ROW,
  BILIN7\_PAIRWISE, POW5 and MODARITH7 (L2296--2750), validator
  tests (L1521--1559) and a check that every opcode has a template
  arm (L1445). They execute a Rust model interpreter
  (\texttt{eval\_\allowbreak{}quotient\_\allowbreak{}vm\_\allowbreak{}for\_\allowbreak{}test}, L3093), so they pin the
  encoder, not the Yul arms; the whole-program test interpreter
  (\texttt{eval\_\allowbreak{}quotient\_\allowbreak{}program\_\allowbreak{}artifacts\_\allowbreak{}for\_\allowbreak{}test}, L3005)
  panics on dynamic opcodes such as MODARITH7 (L3050), so it cannot
  evaluate the production program. There is no EVM-level per-arm test.
\end{itemize}

\paragraph{Fix.}
\begin{itemize}[leftmargin=1.4em,itemsep=2pt]
\item \textbf{Keep the differential harness and wire it into the
  root CI} --- it is the right design, and today nothing gates on it.
\item \textbf{Differentially test every shipped program}: render a
  trace build of each production circuit (B64 included) and run the
  bidirectional trace comparison. Because the VM is data-oblivious,
  one passing run covers every rendered arm and sub-block. A
  render-time census (\texttt{used\_ops}, \texttt{vm/mod.rs:1217},
  plus MODARITH7 flag/block usage) makes it cheap to show which
  fixture covers which artifact. This turns ``we tested some
  shapes'' into ``we tested the program we ship.''
\item \textbf{Run the existing Rust-model opcode tests against the
  Yul arms} (same bytecode, executed in revm), so a Yul-arm bug is
  caught directly rather than only through a fixture.
\item For the residual gap, a \textbf{formal model of the VM
  ISA} (the opcode set is small enough) would close it
  definitively.
\end{itemize}

% ---------------------------------------------------------------------------
\section{C3 --- Truncated challenges: the soundness level is unstated}
\label{sec:c3-soundness}
\textcolor{mathpurple}{\textbf{[MAJOR]}} \vstatus{OPEN}
(arithmetic \vstatus{CORRECTED}: the bit-level figures of earlier drafts are withdrawn)

The \texttt{truncated-challenges} feature masks \texttt{x3} to
128 bits after squeeze (B64 L1270; Rust \texttt{poly/\allowbreak{}kzg/\allowbreak{}mod.rs:457})
and the powers of \texttt{x1}/\texttt{x4} to 128 bits at use
(\texttt{mod.rs:394, 502}); \texttt{x}, \texttt{y} and \texttt{x2}
stay full width. The verifier is faithful to the Rust: the same
truncation happens on both sides. The soundness level this leaves is not established here. The
library's own doc comment (\texttt{proofs/\allowbreak{}src/\allowbreak{}utils/\allowbreak{}arithmetic.rs:242--245})
warns that 128 bits ``may not be enough entropy depending on the
application'' and cites a collision attack with $2^{64}$ memory and
$\sim 2^{64}$ operations. Earlier drafts of this chapter derived a
specific bit level from a Schwartz--Zippel bound; that derivation is
withdrawn, and no replacement figure is given.

\paragraph{The real problem: the silence.}
The critique is about \emph{governance}. The contract ships with
\textbf{no security-level statement}. The only trace in the artifact
is a code comment at L1261--1269 saying the truncation ``mirrors
midnight-proofs''. The header NatSpec, the generator docs and
\texttt{docs/\allowbreak{}audit/\allowbreak{}CODEGEN\_\allowbreak{}ASSURANCE\_\allowbreak{}DOSSIER.md}
state no security level. Truncation also buys no EVM gas: EIP-2537 prices
an MSM independently of scalar width. It exists for in-circuit
verification and is a global crate feature. A tradeoff that is
\textbf{invisible to the consumer} is indistinguishable from an
accident.

\paragraph{Fix.}
\begin{itemize}[leftmargin=1.4em,itemsep=2pt]
\item \textbf{Write down the soundness argument} for the truncated
  multi-open challenges, and state the resulting level in the contract
  header.
\item \textbf{Record a security-team sign-off} that the level is
  acceptable for this deployment's threat model. If it is not, the fix is
  full-width challenges, which cost nothing on-chain for the outermost
  wrap proof but need per-proof configuration in \texttt{midnight-proofs}.
\item \textbf{Surface the truncation in the provenance header}
  (C8) so a deployment audit can enumerate which verifiers are
  truncated.
\end{itemize}

% ---------------------------------------------------------------------------
\section{C4 --- Load-bearing invariants are unenforced}
\label{sec:c4-invariants}
\textcolor{mathpurple}{\textbf{[MAJOR]}} \vstatus{PARTIAL} \emph{(downgraded from
CRITICAL --- the invariant holds today; the gap is that nothing
keeps it holding)}

The contract's safety rests on two invariants:
\begin{itemize}[leftmargin=1.4em,itemsep=2pt]
\item \emph{Every G1ADD operand is a G1MSM output, a VK constant,
  or a proof point that an earlier G1MSM on the same
  \texttt{success} chain has already validated}
  (\S\ref{sec:pr-g1add-subgroup}) --- the thing that makes the
  subgroup-cancellation attack unreachable. (The stricter wording
  ``MSM output / VK constant'' is literally false for L3550: its
  \texttt{PAIRING\_LHS} operand is $\pi$ copied raw at L3471; $\pi$
  is safe because \texttt{G1MSM}$(x_3,\pi)$ at L3486 has already
  subgroup-checked it.) \vstatus{OPEN}: nothing asserts this.
\item \emph{Every transcript-absorbed point flows to a checking
  precompile} --- the thing that makes the deferred transcript
  check sound. \vstatus{ADDRESSED} at codegen time:
  \texttt{validate\_\allowbreak{}generator\_\allowbreak{}invariants} (\texttt{plan.rs:238},
  run at \texttt{plan.rs:164}) calls
  \texttt{kzg::validate\_\allowbreak{}absorbed\_\allowbreak{}g1\_\allowbreak{}precompile\_\allowbreak{}coverage}
  (\texttt{kzg/mod.rs:602}), which fails the render if an absorbed
  proof G1 is missing from the precompile-input set. Caveat: it
  checks the generator's \emph{model} of those inputs, not the
  emitted Yul (the model credits $\pi$ to ``final pairing input'',
  \texttt{kzg/mod.rs:560}, although with the accumulator batch
  $\pi$ reaches the pairing only through the L3550 G1ADD; its
  actual subgroup gate is the L3486 G1MSM).
\end{itemize}

\paragraph{Why this is MAJOR, not CRITICAL (measured).}
These invariants \textbf{hold today} --- the G1ADD cancellation
attack was traced and shown \emph{unreachable}
(\S\ref{sec:pr-g1add-subgroup}): every G1ADD operand is a
point already in $G_1$ (a G1MSM output, or $\pi$ after its G1MSM
check), so no operand carries a cofactor component. There is \textbf{no
current vulnerability}. A CRITICAL rating requires a live flaw;
this is a \emph{defense-in-depth gap} --- the safety is real but
unprotected. The honest grade is MAJOR.

\paragraph{The real risk: silent future breakage.}
The danger is not today, it is the \textbf{next codegen
refactor}. A future change that emits a G1ADD with a raw proof
point that no earlier G1MSM has validated turns a mathematically-possible attack into a
live one, and \textbf{no test or assertion catches it} because
the G1ADD half of the invariant was never written down as a check. The safety is a
property of the current emission, not an enforced contract on
future emission. That is fragile, and fragility in a soundness
boundary is worth fixing now.

\paragraph{Fix.}
\textbf{Codegen-time assertions that fail the build.} The absorbed-point
half already exists (above); add the G1ADD half: the
generator should statically verify, before rendering, that every
G1ADD operand is an MSM output/constant or a point validated by an
earlier MSM on the same success chain, and derive both checks from
the emitted precompile calls rather than a parallel model. An invariant that
isn't machine-checked is a comment, not a guarantee. This is
cheap (a static pass over the emitted precompile graph) and converts
``it happens to be safe'' into ``it cannot be emitted unsafe.''

% ---------------------------------------------------------------------------
\section{C5 --- The contract defeats stock automated tools}
\label{sec:c5-tools}
\textcolor{mathpurple}{\textbf{[MAJOR]}} \vstatus{PARTIAL}

Stock slither 0.11.6 \emph{crashes} in the IR$\to$SSA pass
(\texttt{RecursionError} in \texttt{add\_phi\_origins}); aderyn
false-positives on assembly patterns (its own L-3 note admits it
cannot see assembly references); semgrep is blind. A contract
that \textbf{stock linters cannot analyze} is a contract that can only
be audited by expensive manual review --- which does not scale
and is error-prone. ``Semgrep clean'' here is not assurance; it
is the absence of a capable observer. The assembly-VM choice
directly causes this.

\paragraph{Correction (measured).} The crash is not caused by
switch nesting. SSA construction recurses once per
dominator-tree level, and a straight-line Yul body is a dominator
\emph{chain}; \texttt{verifyProof} has 2{,}816 CFG nodes. With
three small patches (iterative dominator walks, a deep-recursion
worker thread, Yul \texttt{leave} CFG edges), slither analyses B64
completely (\texttt{nightfall\_\allowbreak{}4\_\allowbreak{}PV/\allowbreak{}audit/\allowbreak{}slither-2026-09-29},
\texttt{runs/\allowbreak{}generated\_\allowbreak{}B64\_\allowbreak{}final.txt}). Its
\texttt{unused-state} detector then resolves the assembly
references correctly and flags exactly the 5 pointer constants
that really are unreferenced (\texttt{CHALLENGE\_MPTR},
\texttt{Q\_COM\_MPTR}, \texttt{G1\_\allowbreak{}IDENTITY\_\allowbreak{}MPTR},
\texttt{QUOTIENT\_\allowbreak{}RETURN\_\allowbreak{}MPTR}, \texttt{TRACE\_U256\_MPTR}); it
also reports a dead store (L2003--2004), and the audit's
Slither-API pass confirms that all 18 \texttt{staticcall}s check
\texttt{returndatasize}. So aderyn's 76 L-3 hits are 71 false
positives and 5 true ones, not all false. The residual critique
stands: the tools see Yul arithmetic but nothing of the proof
system, so ``tool-clean'' is still not assurance.

\paragraph{Fix.}
\begin{itemize}[leftmargin=1.4em,itemsep=2pt]
\item Upstream the three slither patches, and shorten the
  dominator chain at the source: splitting \texttt{verifyProof}
  into functions (R3 / \S\ref{sec:read-decompose}) bounds each
  CFG. Splitting only the VM \texttt{switch} would not help.
\item Ship a \textbf{Solidity-level reference implementation}
  alongside the assembly, used for tooling and differential
  testing against the assembly hot loop.
\item Do not let ``tools can't run'' become the accepted state ---
  a contract that resists all tooling is a contract that resists
  maintenance.
\end{itemize}

% ---------------------------------------------------------------------------
\section{C6 --- Compiler-internal names leak into the audited artifact}
\label{sec:c6-names}
\textcolor{mathpurple}{\textbf{[MAJOR]}} \vstatus{OPEN}

The contract carries \textbf{163 distinct \texttt{varN} names}
(980 uses; 442 \texttt{let varN} declarations, because numbering
restarts per emitted block and the same name denotes different
values in different blocks) and
\textbf{31 \texttt{f\_N}/\texttt{a\_N} names} --- generator-internal
identifiers emitted into the file a human must read. (The
\texttt{f\_N}/\texttt{a\_N} names are the one place the generator
already passes a semantic hint: fixed/advice column $N$,
\texttt{yul\_emit.rs:1142}.) The
generator \emph{knows} the semantic role of every variable (it
assigned it) and still emits \texttt{var47}. This is the
generator being lazy about the audit surface: it optimizes its own
emission simplicity over the reviewer's comprehension.

\paragraph{Fix.}
Name variables semantically in the emitted code --- the generator
has the mapping (for gate subexpressions, the gate and constraint
names of the quotient identity manifest). \texttt{var47} $\to$
\texttt{perm\_\allowbreak{}z\_\allowbreak{}numerator\_\allowbreak{}term}. This is a pure codegen change
with zero runtime cost and a large readability win (\S\ref{sec:read-varnames}). The same
applies to the anonymous VM scratch addresses
(\S\ref{sec:pr-legend}, \S\ref{sec:read-named-addr}).

% ---------------------------------------------------------------------------
\section{C7 --- The pairing helper's parameter names invert the slot names}
\label{sec:c7-naming}
\textcolor{codegenblue}{\textbf{[MINOR]}} \vstatus{CORRECTED}

\texttt{ec\_\allowbreak{}pairing(success, PAIRING\_\allowbreak{}RHS, PAIRING\_\allowbreak{}LHS)} (L3569)
passes the slots in the opposite order to the helper's parameter
names \texttt{lhs\_\allowbreak{}mptr, rhs\_\allowbreak{}mptr} (L534). The original critique
called the slot names ``backwards'' and the rename ``never made''.
That is wrong on both counts. The slot names follow the Rust
\texttt{DualMSM}: \texttt{left} $=\pi$, \texttt{right} $=$
\texttt{final\_com} $- vG + x_3\pi$ (\texttt{poly/\allowbreak{}kzg/\allowbreak{}mod.rs:554--558}),
and the differential trace pins them as ids 27/28
``pairing\_lhs''/``pairing\_rhs''. Rust checks
$e(\text{left},[s]_2)\,e(\text{right},-[1]_2)=1$ (\texttt{msm.rs:284--300});
Solidity checks the equivalent $e(\text{right},[1]_2)\,e(\text{left},-[s]_2)=1$.
The swap is explained at the call site in a 12-line comment (L3555--3567,
template \texttt{FinalPairing.yul:65--79}). A wrong argument order
would reject every honest proof, which every positive test catches,
so it is not a latent soundness hazard. Renaming the slots
would break the name-level correspondence with Rust and the trace.

\paragraph{Fix.}
Rename the \emph{helper's parameters}, not the slots:
\texttt{ec\_\allowbreak{}pairing(success, g2\_\allowbreak{}one\_\allowbreak{}side\_\allowbreak{}mptr, neg\_\allowbreak{}s\_\allowbreak{}side\_\allowbreak{}mptr)}
(\texttt{AssemblyHelpers.yul:213}) removes the trap without breaking
the correspondence. If semantic slot names are wanted, use
\texttt{KZG\_PI} / \texttt{KZG\_\allowbreak{}COMBINED\_\allowbreak{}COMMITMENT}. The earlier
suggestion \texttt{KZG\_EVAL\_POINT} was wrong: $\pi$ is the opening
witness commitment, not an evaluation point.

% ---------------------------------------------------------------------------
\section{C8 --- No provenance metadata}
\label{sec:c8-provenance}
\textcolor{mathpurple}{\textbf{[MAJOR]}} \vstatus{OPEN}

The contract records \textbf{no generator version, no Rust commit
hash, no generation timestamp}. It does carry content-derived
identifiers: the VK's \texttt{vk\_digest} (the transcript
representation of the constraint system), \texttt{EXPECTED\_\allowbreak{}VK\_\allowbreak{}CODEHASH\_\allowbreak{}WORD}
(L46), and, in the B64 lineage, \texttt{EXPECTED\_\allowbreak{}QUOTIENT\_\allowbreak{}CERTIFICATE\_\allowbreak{}HASH\_\allowbreak{}WORD}
(L51). These identify \emph{which circuit/VK} a verifier accepts,
but not \emph{which generator build} produced it. For a generated
artifact this is unmaintainable: you cannot tell whether a given
deployed contract was produced by the current generator or a
stale one. When a bug is found in the generator, you have
no way to enumerate the affected deployments. The project's own
\texttt{docs/\allowbreak{}audit/\allowbreak{}CODEGEN\_\allowbreak{}ASSURANCE\_\allowbreak{}DOSSIER.md} already lists the
record a production claim needs (Midfall revision, toolchain, Cargo
features, solc version and flags); the artifact just does not carry it.

\paragraph{Fix.}
Embed a provenance header as constants plus a comment:
\texttt{GENERATOR\_\allowbreak{}VERSION}, \texttt{RUST\_COMMIT},
\texttt{CARGO\_FEATURES} (e.g.\ \texttt{truncated-challenges}),
\texttt{GENERATED\_AT}, and name the existing
\texttt{vk\_digest}/codehash words as the circuit id. Unreferenced
constants do not change the runtime bytecode, but any source edit
changes the CBOR metadata hash unless metadata is omitted (the
generator's builds pass \texttt{--no-cbor-metadata},
\texttt{src/evm.rs:209}); keep volatile fields such as the timestamp
in a sidecar manifest. Make the
generator stamp every artifact. This is table stakes for
generated code.

% ---------------------------------------------------------------------------
\section{C9 --- Magic constants hand-synced with no assertion}
\label{sec:c9-magic}
\textcolor{codegenblue}{\textbf{[MINOR]}} \vstatus{ADDRESSED}

\emph{Original claim:} \texttt{EXPECTED\_\allowbreak{}VK\_\allowbreak{}PAYLOAD\_\allowbreak{}LENGTH = 17056}, the permutation
$\delta$, \texttt{q\_\allowbreak{}perm\_\allowbreak{}chunk\_\allowbreak{}len} must be kept in
sync with the VK by the generator, but nothing in the
contract asserts they match the loaded VK.

\emph{Re-verified:} the length \emph{is} asserted.
\texttt{EXPECTED\_\allowbreak{}VK\_\allowbreak{}LENGTH} $=$ payload $+1$ (the \texttt{INVALID}
prefix) is checked together with the full-runtime codehash at
construction (L266--267) and again on every proof (L921--922,
template \texttt{VkLoading.yul:67--70}); the codehash covers every
payload byte. After the copy, the VK header words
\texttt{NUM\_INSTANCES}, \texttt{K}, \texttt{HAS\_ACCUMULATOR},
\texttt{ACC\_OFFSET}, \texttt{NUM\_ACC\_LIMBS} and
\texttt{NUM\_\allowbreak{}ACC\_\allowbreak{}LIMB\_\allowbreak{}BITS} are cross-checked against the
generator's constants (L934--939, \texttt{VkLoading.yul:82--88},
commit \texttt{e64f2891}); the B64 lineage also checks the
quotient-certificate hash (L940). $\delta$ is not VK data: it is the
field constant \texttt{Fq::DELTA} (\texttt{quotient.rs:1379}).
\texttt{q\_\allowbreak{}perm\_\allowbreak{}chunk\_\allowbreak{}len} is shape metadata from the same
constraint system that the pinned VK digests.

\paragraph{Fix (residual).}
None needed for soundness. Optionally derive
\texttt{q\_\allowbreak{}perm\_\allowbreak{}chunk\_\allowbreak{}len} from a VK header word and cross-check it
like the others.

% ---------------------------------------------------------------------------
\section{C10 --- The negated-numerator sign convention is a footgun}
\label{sec:c10-sign}
\textcolor{codegenblue}{\textbf{[MINOR]}} \vstatus{CORRECTED}

The verifier stores $-\nu_y(x)$ at \texttt{QUOTIENT\_\allowbreak{}EVAL\_\allowbreak{}MPTR}
(L2823--2824), not at \texttt{QUOTIENT\_MPTR}, which in this path
holds \texttt{x\_split} and \texttt{one\_minus\_x\_n} (L2865--2866).
The trace re-negates it for id 36 (\texttt{TraceReturn.yul:37}). The
original critique said a dropped or doubled negation would
``silently'' accept wrong proofs or reject right ones. It would not
be silent. The negation is not a Solidity trick: it is Rust's
\texttt{expected\_eval}, into which
\texttt{compute\_\allowbreak{}linearization\_\allowbreak{}commitment} subtracts the fully
evaluated identities. The value is pinned by the differential
trace (id 36 against \texttt{verifier.rs:508--524}), and every
positive end-to-end test fails if the sign flips. The convention is
documented three times (L120--122, L1460--1461, L2820--2822). What
remains is a naming problem: the constant's own comment calls
\texttt{QUOTIENT\_\allowbreak{}EVAL\_\allowbreak{}MPTR} a ``legacy name''.

\paragraph{Fix.}
Rename the slot after what it holds
(\texttt{LINEARIZATION\_\allowbreak{}EXPECTED\_\allowbreak{}EVAL\_\allowbreak{}MPTR}) and drop the legacy
\texttt{QUOTIENT\_MPTR} name in the fused path. Keep the negation,
since it matches the Rust oracle term for term.

% ---------------------------------------------------------------------------
\section{C11 --- Build fragility: \texttt{via-ir} is load-bearing}
\label{sec:c11-build}
\textcolor{mathpurple}{\textbf{[MAJOR]}} \vstatus{PARTIAL}
\emph{(upgraded from MINOR: see the spill-area paragraph)}

The contract \textbf{will not compile with the default codegen}
(re-verified with solc 0.8.30 without \texttt{--via-ir}:
``Stack too deep \ldots\ Variable \texttt{src} is 1 slot(s) too
deep''). A compiler flag is load-bearing for the build.
If CI misconfigures, the contract fails to build --- which is
fail-closed (good) --- but it signals the code is operating at
the edge of the compiler's limits.

\paragraph{Already pinned} \vstatus{ADDRESSED}. The generator
compiles only with \texttt{--via-ir} (\texttt{src/evm.rs:206});
the flag set is recorded in
\texttt{docs/\allowbreak{}audit/\allowbreak{}REVIEW\_\allowbreak{}PACKET.md:71} and the assurance dossier;
Nightfall's \texttt{foundry.toml} (L31--50) pins \texttt{via\_ir}
per file for the B8/B64 verifiers through
\texttt{compilation\_\allowbreak{}restrictions}, citing the stack-too-deep error.

\paragraph{The spill area (new, measured).} Under
\texttt{--via-ir} (solc 0.8.30 and 0.8.35, 200 runs) the deployed
code starts with
\texttt{PUSH2 0x08e0 PUSH1 0x40 MSTORE}: solc's stack-limit evader
has reserved the 67-word area $[\texttt{0x80},\texttt{0x8e0})$ for Yul
locals it moved to memory (for example \texttt{y} at \texttt{0x300} and the
MODARITH7 pairwise inner-loop counter at \texttt{0x80}). The generated code writes
the same range as \texttt{TRANSCRIPT\_MPTR} $=$
\texttt{RETURN\_MPTR} $=$ \texttt{0x80}, the accumulator batch at
\texttt{0x100}--\texttt{0x320}, and the \texttt{ec\_pairing} scratch
at \texttt{0x300}--\texttt{0x600}. So the
\texttt{assembly ("memory-safe")} promise, which the generator makes
precisely so via-IR can spill (\texttt{tests.rs:2127}), is formally
broken. It works today only because the spilled locals are live
in the quotient phase alone, where no generated write goes below
\texttt{0x8e0}. That is consistent with every fixture passing, but
nothing checks it, and any Yul restructuring (e.g.\ R3) or solc
upgrade reshuffles the spill set.

\paragraph{Fix.}
Keep \texttt{via-ir}; removing it is not realistic. The first legacy
error is in the 6-parameter helper \texttt{load\_\allowbreak{}acc\_\allowbreak{}coord\_\allowbreak{}shifted}
(L573), not in the VM dispatch, and the gas tuning assumes via-IR.
Instead, close the spill hazard. Either
start generated absolute scratch above a fixed margin and add a CI
gate that parses the deployed \texttt{memoryguard} and fails if it
exceeds that base, or make the planner reserve
$[\texttt{0x80},\textit{memoryguard})$ from a first compilation pass.

% ---------------------------------------------------------------------------
\section{C12 --- ``Generated'' hides that the real artifact is the compiler}
\label{sec:c12-framing}
\textcolor{mathpurple}{\textbf{[MAJOR]}} \vstatus{PARTIAL}

The framing ``we generate the verifier'' obscures where the
complexity actually lives. The contract is small and opaque; the
\textbf{semantic content is in the generator} --- the shape
recognizer, the fusion decision, the run compactor, the offline
validator (\S\ref{sec:pr-offline-validator} onward). Auditing the
contract does not audit the system; \textbf{auditing the
generator does}. The VM's auditability win is an illusion: the
complexity was not removed, it was \textbf{relocated to a
harder-to-audit place}. (The earlier claim that the generator gets
the \emph{least} review discipline was too strong. It has
$\approx$187 \texttt{\#[test]} functions, a pre-render invariant
pass (\texttt{plan.rs:238}), proptest-based property tests, a
pinned runtime hash for the IVC fixture
(\texttt{scripts/\allowbreak{}check\_\allowbreak{}release\_\allowbreak{}bytecode\_\allowbreak{}sizes.sh}), and
\texttt{docs/\allowbreak{}audit/\allowbreak{}CODEGEN\_\allowbreak{}ASSURANCE\_\allowbreak{}DOSSIER.md}. What it lacks is
a CI that runs any of this on PRs (\S\ref{sec:c2-trustroot}) and
the items below.)

\paragraph{Fix.}
Apply the same review rigor to the generator that you would to a
hand-written verifier. Property-based tests over the lowering
passes exist in part; still missing are a golden-program corpus
with expected disassembly (one pinned hash per shipped artifact,
not only the IVC fixture) and a formal model of the VM ISA. Treat
the generator as the product and the contract as its output
artifact.

% ---------------------------------------------------------------------------
\section{What the design actually gets right}
\label{sec:c-balance}

To be fair --- and a critique that only lists faults is not
credible --- the design has genuine strengths that should be
preserved:
\begin{itemize}[leftmargin=1.4em,itemsep=2pt]
\item \textbf{The differential trace-equivalence harness} is the
  standout. The crate's own workflow file defines a PBT differential
  fuzzer and dedicated native-vs-Solidity trace-equivalence jobs
  (Poseidon, IVC) with bidirectional id-and-value matching.
  This is more assurance than most verifier projects ship, and it
  is the reason the ``correct'' verdict in this review is
  trustworthy --- once it actually gates PRs (it is not wired into
  the monorepo's root CI; \S\ref{sec:c2-trustroot}).
\item \textbf{The fail-closed dispatch chain}
  (\S\ref{sec:pr-failclosed}) is exemplary: unknown opcodes,
  unknown sub-cases, and non-terminal VM states all revert. This
  is the right default and it is consistently applied.
\item \textbf{The double codehash anchor}
  (\S\ref{sec:pr-vk-generic}) --- checked at both the Solidity
  and Yul levels --- is a clean, defense-in-depth trust binding.
\item \textbf{The offline bytecode validator}
  (\S\ref{sec:pr-offline-validator}) --- a static analyzer that
  checks stack effects before rendering --- is a genuinely good
  idea that most codegen projects lack. The same holds for the
  pre-render invariant pass (\texttt{plan.rs:238}: memory
  non-overlap, absorbed-G1 precompile coverage, PCS/quotient
  drift), which fails the build rather than emitting.
\item \textbf{The constant interning and run compaction} are real
  engineering, not hand-waving.
\end{itemize}
These show the team \emph{can} do rigorous work. The critique
above is that the rigor is applied unevenly: the generator's
internal machinery and its differential harness are strong, while the
artifact it produces is left opaque, the trust root is
differentially tested but not formally verified, and the
produced contract resists stock tooling.

% ---------------------------------------------------------------------------
\section{Measured auditability scorecard}
\label{sec:c-scorecard}

To keep the critique honest, the readability findings are
quantified against the B64 artifact (3{,}579 lines). Numbers,
not vibes. Counts are over code with \texttt{//} comments stripped.
The memory-operand metric is the first argument of
\texttt{mload}/\texttt{mstore}/\texttt{mcopy} (914 such ops) when
that argument is a bare constant: 492 raw \texttt{0x}\ldots\
literals (246 distinct addresses) against 64 named
\texttt{*\_MPTR}/\texttt{*\_CPTR}. Counting \texttt{add(BASE, off)}
forms too gives 514 raw against 176 named. The same metric is used
in \S\ref{sec:read-measure}.

\begin{center}\small
\begin{tabular}{@{}p{5.6cm}rp{5.2cm}@{}}
\toprule
\textbf{Metric} & \textbf{Value} & \textbf{Reading} \\
\midrule
Comment : code lines & 47.2\% & \textcolor{tracegreen}{\textbf{Good}} --- 1{,}109 comment vs 2{,}350 code lines (120 blank) \\
Named constants & 79 & \textcolor{tracegreen}{\textbf{Good}} --- 68 pointer constants (\texttt{*\_MPTR}/\texttt{*\_CPTR}/\texttt{*\_MPTR\_BASE}), 63 referenced, 5 dead \\
\texttt{varN} SSA names & 163 distinct / 980 uses & \textcolor{rustorange}{\textbf{Bad}} --- generator names leak (C6) \\
Raw-hex literals & 1{,}409 (573 distinct) & \textcolor{rustorange}{\textbf{Bad}} --- 0.60 per code line \\
Memory-operand naming & 11.5\% named / 88.5\% raw hex & \textcolor{rustorange}{\textbf{Bad}} --- 64 named vs 492 raw (246 distinct) (I9) \\
Memory ops & 914 (480 / 336 / 98) & \texttt{mload} / \texttt{mstore} / \texttt{mcopy}; scratch traffic is the bulk of the code \\
\texttt{verifyProof} length & \textbf{3{,}286 lines} & \textcolor{rustorange}{\textbf{Bad}} --- god-function (C1) \\
Nested Yul functions & 17 & helpers are named (mitigates C6) \\
Branch points & 2 switch / 15 case / 73 if & VM dispatch + guards \\
\bottomrule
\end{tabular}
\end{center}

\paragraph{What the numbers say.}
The contract is \emph{heavily commented} (47\% --- genuinely
helpful) and the \emph{high-level} memory map is named (68
pointer constants). But the \emph{working} code is dominated by
anonymous scratch: \textbf{88.5\% of bare-constant memory
operands are raw hex}, 163 \texttt{varN} names carry the
dataflow, and the whole verifier is one \textbf{3{,}286-line
\texttt{verifyProof}}. The comments describe the code; they do
not substitute for naming it. A reviewer reads 1{,}109 lines of
prose and then still has to decode 492 raw-hex memory operands
(246 distinct addresses: 183 in the final-MSM/quotient scratch
from \texttt{0xa320}, 60 in the decoded-eval buffer from
\texttt{0x8520}, 3 low) and 980 \texttt{varN} references by hand.

\paragraph{The confirmed C2 coverage gap (measured, corrected).}
The VM ISA defines \textbf{31 opcodes} (\texttt{Q\_\allowbreak{}OP\_\allowbreak{}PUSH\_\allowbreak{}CONST}
\ldots \texttt{Q\_\allowbreak{}OP\_\allowbreak{}AFFINE\_\allowbreak{}SUM}, \texttt{vm/\allowbreak{}mod.rs:560--590}), but a
render contains only the arms its program uses. B64 has 11
(\texttt{0x05, 0x06, 0x08, 0x0b, 0x0d, 0x10, 0x11, 0x19, 0x1b,
0x1f, 0x21}); \texttt{BILIN7\_PAIRWISE}, \texttt{POW5} and the
other 18 are not rendered and would hit the default revert (L2561),
so they cannot be ``untested code in production''. B64's MODARITH7
instructions use the mem/product (``escape-hatch'') blocks heavily
and never the pairwise block. The differential test's
required-coverage check (\texttt{src/\allowbreak{}test.rs:1697--1743}) asserts
\emph{semantic verifier stages} --- challenge ids 7--11 and 13--16,
quotient numerator (36), f\_eval (31),
final MSM (33), pairing lhs/rhs (27/28), final pairing (35),
and trace ranges 30k--42k / 60k--61k --- and compares every
per-identity value, but it records \textbf{no arm or sub-block
coverage}. Since the VM is data-oblivious, that is enough for any
program that is tested at all. The closest fixture, the IVC Keccak
decider, renders the same 11-opcode histogram as B64, so B64's arms
are exercised \emph{in shape}; B64's exact program, $k{=}21$ layout and
embedded-quotient mode are not. None of it is gated by the monorepo's
root CI. This is the precise, measurable form of C2's residual:
\emph{a tested shape $\neq$ the tested shipped program}.

\paragraph{Composite verdict.}
Readability is \textbf{Fair}: strong comments and a named
high-level map, undermined by anonymous scratch (88.5\% raw-hex
bare memory operands), leaked generator names, a 3{,}286-line
god-function, and a differential test that checks stages but not
rendered arms. Every one of these is a codegen-emission choice, not a
runtime necessity --- the same bytecode could be emitted with
named slots, semantic variable names, split functions, and an
arm-coverage assertion (split functions would need the spill-area
re-check of C11). The gap between ``correct'' and
``auditable'' is entirely a matter of what the generator chooses
to emit.

\section{Audit runbook: how to actually audit this contract}
\label{sec:c-runbook}

The critique says the contract is hard to audit. The constructive
capstone is a concrete procedure that makes it auditable
\emph{today}, working around the tooling gaps with the
assurance the project already provides. A reviewer with a day
and this list can reach a sound verdict.

\begin{enumerate}[leftmargin=1.6em,itemsep=4pt]
\item \textbf{Build integrity.} Compile with
  \texttt{solc --via-ir --optimize} (the contract will not build
  without \texttt{via-ir} --- stack-too-deep). Confirm the
  runtime is \texttt{<24\,KB} (measured 21{,}253\,B with CBOR
  metadata, 21{,}200\,B without, solc 0.8.30, 200 runs). The
  pragma is \emph{not} pinned (\texttt{\^{}0.8.24}), so record the
  exact compiler (generator: 0.8.30, \texttt{src/evm.rs:78};
  Nightfall: 0.8.35, \texttt{foundry.toml}). Read the deployed
  prefix \texttt{PUSH2 }$g$\texttt{ PUSH1 0x40 MSTORE} and confirm
  no spilled local is live across generated writes below $g$
  (C11; $g=$\texttt{0x8e0} today). \emph{Evidence:} \texttt{solc} output,
  \texttt{--bin-runtime} size and prefix.
\item \textbf{VK binding (the trust anchor).} Verify
  \texttt{EXPECTED\_\allowbreak{}VK\_\allowbreak{}CODEHASH\_\allowbreak{}WORD} matches the deployed
  VK's \texttt{extcodehash} and length, at \emph{both} the Solidity level
  (L266--267) and the Yul level (L921--922), and that the header
  cross-checks (L934--940) match the VK words. A mismatch here breaks
  everything downstream. \emph{Evidence:}
  \texttt{getCode(vk)} vs the constant.
\item \textbf{Read the named skeleton.} Walk the 68 named
  pointer constants (63 referenced; VK header, challenges
  \texttt{THETA/\allowbreak{}BETA/\allowbreak{}GAMMA/\allowbreak{}Y/\allowbreak{}X}, pairing inputs). This is the
  readable high-level map --- understand the data flow here
  before touching the VM.
\item \textbf{Transcript schedule cross-check.} Compare the
  Fiat--Shamir absorb order against the Rust verifier's
  schedule (\S\ref{sec:rust-transcript} / \S\ref{sec:transcript-schedule}). The order is
  soundness-critical; a reordering is silent. \emph{Evidence:}
  the differential trace (step 7) also covers this.
\item \textbf{Fail-closed dispatch.} Confirm the three-layer
  revert chain: main opcode switch default (L2561), heavy-gate
  u16 sub-dispatch default (L2528), and the terminal guards
  (\texttt{q\_pc==q\_end} and \texttt{!q\_has\_top}, L2569--2570).
  Check that the 11 rendered arms match the decoded program's
  opcode census.
  \emph{Evidence:} \S\ref{sec:pr-failclosed}.
\item \textbf{G1ADD operand provenance.} For each of the 4
  runtime G1ADDs (L3480, L3490, L3536, L3550), confirm each operand is a
  G1MSM output, a VK constant, or a proof point already validated by an
  earlier G1MSM on the same \texttt{success} chain (L3550's
  \texttt{PAIRING\_LHS} is raw $\pi$, validated by the L3486 G1MSM;
  \S\ref{sec:pr-g1add-subgroup}, \S\ref{sec:c4-invariants}). This is the one manual check
  that must not be skipped --- it is what makes the
  subgroup-cancellation attack unreachable.
\item \textbf{Run the differential harness.} Execute the
  native-vs-Solidity trace-equivalence jobs (Poseidon + IVC)
  with \texttt{rust-verifier-trace,solidity-trace} and
  \texttt{HALO2\_\allowbreak{}SOLIDITY\_\allowbreak{}RUN\_\allowbreak{}EVM\_\allowbreak{}TESTS=1}. This is the
  strongest single assurance --- bidirectional trace matching
  against the Rust oracle. \emph{Evidence:} the job definitions in
  \texttt{proofs/\allowbreak{}solidity-verifier/\allowbreak{}.github/\allowbreak{}workflows/\allowbreak{}ci.yml}, which
  GitHub does not execute in this monorepo, so run them yourself. They
  validate the \emph{generator} on its fixtures; for B64 itself,
  render a trace build of the fee circuit and run the same
  comparison, since nothing else does.
\item \textbf{VM scratch legend.} Use the named-constant legend
  (\S\ref{sec:pr-legend}) or the R1 disassembly to map the
  $\approx$45 distinct raw-hex addresses in the VM body (L1692--2564)
  to witness/permutation columns.
  Without this, the VM body is unreadable.
\item \textbf{Tooling triage.} Run aderyn/semgrep and clear the
  \emph{known} false positives: shift-order (H-1),
  return-in-Yul (H-2), unused-state-var (L-3, except the 5 genuinely
  dead constants of C5). Run slither with the three patches of
  \texttt{nightfall\_\allowbreak{}4\_\allowbreak{}PV/\allowbreak{}audit/\allowbreak{}slither-2026-09-29}. Do not treat
  ``semgrep clean'' as assurance --- it cannot see the assembly
  (\S\ref{sec:pr-tooling}).
\item \textbf{Security level.} Confirm the truncation
  (\S\ref{sec:c3-soundness}) has a written soundness argument, is
  labeled, and is signed off for the deployment's threat model.
\end{enumerate}

\paragraph{Why this runbook works.}
It substitutes the project's \emph{own} strong assurance (the
differential harness, the codehash anchor, the fail-closed chain)
for the tooling that cannot reach the assembly. The two checks
that genuinely require human attention are \#2 (the trust
anchor) and \#6 (the G1ADD provenance) --- everything else is
either mechanical or covered by the harness. That is a tractable
audit surface for a contract that stock linters cannot analyze.

\section{Critique summary}
\label{sec:c-summary}

\begin{center}\small
\begin{tabular}{@{}p{0.27\linewidth}ccp{0.40\linewidth}@{}}
\toprule
\textbf{Critique} & \textbf{Sev} & \textbf{Status} & \textbf{Concrete fix} \\
\midrule
C1 VM is the wrong abstraction & \textcolor{mathpurple}{MAJOR} & \vstatus{PARTIAL} & Annotated disassembly (R1) + arm-coverage assertion (R2); split \texttt{verifyProof} (R3). Runtime-cost claim corrected (101 instr., dispatch $\lesssim$2\%) \\
C2 Trust root diff-tested, not formally verified & \textcolor{mathpurple}{MAJOR} & \vstatus{PARTIAL} & Wire the harness into root CI; differentially test each shipped program; run the Rust-model opcode tests against the Yul arms \\
C3 Truncated challenges, level unstated & \textcolor{mathpurple}{MAJOR} & \vstatus{OPEN} & Write the soundness argument; state the level in the header + security sign-off \\
C4 Invariants unenforced (hold today) & \textcolor{mathpurple}{MAJOR} & \vstatus{PARTIAL} & Absorbed-point check exists (\texttt{plan.rs:245}); add the G1ADD-provenance check \\
C5 Defeats stock tooling & \textcolor{mathpurple}{MAJOR} & \vstatus{PARTIAL} & Upstream the slither patches; split \texttt{verifyProof} (not the switch); Solidity reference impl \\
C6 Compiler names leak & \textcolor{mathpurple}{MAJOR} & \vstatus{OPEN} & Semantic variable naming in emitted code \\
C7 Pairing helper parameter names & \textcolor{codegenblue}{MINOR} & \vstatus{CORRECTED} & Slot names follow Rust \texttt{DualMSM}; rename the helper's parameters only \\
C8 No provenance & \textcolor{mathpurple}{MAJOR} & \vstatus{OPEN} & Stamp generator version, commit, features; sidecar manifest per the dossier \\
C9 Magic constants unasserted & \textcolor{codegenblue}{MINOR} & \vstatus{ADDRESSED} & Length + codehash checked (L266, L921); header words cross-checked (L934--940) \\
C10 Sign-convention footgun & \textcolor{codegenblue}{MINOR} & \vstatus{CORRECTED} & Mirrors Rust \texttt{expected\_eval}, pinned by trace id 36; rename \texttt{QUOTIENT\_\allowbreak{}EVAL\_\allowbreak{}MPTR} \\
C11 via-ir load-bearing & \textcolor{mathpurple}{MAJOR} & \vstatus{PARTIAL} & via-IR already pinned; gate the \texttt{memoryguard} spill area against generated scratch \\
C12 ``Generated'' hides the real artifact & \textcolor{mathpurple}{MAJOR} & \vstatus{PARTIAL} & Golden per-artifact corpus + formal ISA model on top of the existing tests \\
\bottomrule
\end{tabular}
\end{center}

\paragraph{Bottom line.}
After measurement, \textbf{no critique survived as CRITICAL} ---
and that is the most important finding of this chapter. The
verifier is \textbf{correct} (the generator's differential harness
passes against the Rust oracle, though no CI gates on it; B64 passes
Nightfall's own fixtures), and its safety invariants
\textbf{hold} (the G1ADD cancellation is unreachable). The
design's problems are all \textbf{MAJOR auditability and
governance concerns}, not soundness breaks: an opaque program
blob, a 3{,}286-line god-function, leaked generator names, an
unlabeled security level, a G1ADD invariant that holds by
emission rather than by assertion (the absorbed-point half is now
asserted), and a via-IR spill area that overlaps generated scratch
with nothing checking the live ranges. None of these lose funds
today. All of them make the next bug easier to hide and the
security posture harder to verify.

This is a favorable result delivered harshly: the team built a
\emph{correct} verifier, and the fail-closed chain, the offline
validator, and the differential harness prove they \emph{can} do
rigorous work. The gap is one of \textbf{priority and
communication} --- the rigor went into the generator's internals
but not into making the produced artifact legible, the security
level loud, or the safety invariants enforced. Almost every fix in this
chapter changes only what the generator \emph{emits}, never what
it computes (the C11 spill fix may move scratch addresses). Fix the priority, and the design becomes as
auditable as it is correct.
